\documentclass[12pt]{article}

\usepackage{amsmath, amssymb, amsthm}
\usepackage{geometry}
\usepackage{hyperref}
\usepackage{natbib}
\usepackage{enumitem}
\usepackage{mathtools}
\usepackage{mathrsfs}

\geometry{margin=1.25in}

\newtheorem{theorem}{Theorem}
\newtheorem{definition}{Definition}
\newtheorem{corollary}{Corollary}
\newtheorem{remark}{Remark}

\title{Blackwell's Theorem on the Comparison of Experiments:\\
       A Bayesian Exposition}
\author{Analytical Summary of \citet{kihlstrom1984}}
\date{}

\begin{document}

\maketitle
\tableofcontents
\bigskip

%----------------------------------------------------------------------
\section{Introduction and Motivation}
%----------------------------------------------------------------------

Consider two random variables, $\tilde{x}_\mu$ and $\tilde{x}_\nu$, each correlated
with an unknown state $\tilde{s}$.  A decision maker wishes to know which observation
conveys more information about $\tilde{s}$.  \citet{blackwell1951,blackwell1953}
identified a clean answer: $\tilde{x}_\mu$ is at least as informative as
$\tilde{x}_\nu$ if and only if the conditional distribution of $\tilde{s}$ given
$\tilde{x}_\mu$ is a \emph{mean-preserving spread} of the conditional distribution
of $\tilde{s}$ given $\tilde{x}_\nu$---or, equivalently, if $\tilde{x}_\mu$ is
\emph{sufficient} for $\tilde{x}_\nu$.

\citet{kihlstrom1984} provides a Bayesian restatement of Blackwell's result and
its equivalence with the economic criterion of \citet{bonnenblust1949}, the
sufficiency criterion of Blackwell, and the uncertainty-function criterion of
\citet{degroot1962}.  The exposition proceeds in three steps: (i)~set up notation
and define the experiments; (ii)~state the three equivalent criteria in a unified
Bayesian language; (iii)~prove their equivalence by establishing the central role
of \emph{stochastic transformations} (Markov kernels).

%----------------------------------------------------------------------
\section{Setup and Notation}
%----------------------------------------------------------------------

\subsection{Experiments and Markov matrices}

Let $S = \{s_1, \ldots, s_N\}$ be the finite set of possible values of the
unknown state $\tilde{s}$.  An \textbf{experiment} is described by the
conditional probability distribution of an observed signal $\tilde{x}$ over
a measurable space $(X, \mathscr{F})$, given the state.  When $X$ is also
finite, say $X = \{x_1, \ldots, x_M\}$, an experiment reduces to an
$N \times M$ \textbf{Markov matrix}
\[
\mu = [\mu_{ij}], \qquad
\mu_{ij} = \Pr(\tilde{x}_\mu = x_j \mid \tilde{s} = s_i) \geq 0,
\quad \sum_{j=1}^{M} \mu_{ij} = 1 \;\forall\, i.
\]
Two experiments $\mu$ (the ``more informative'' candidate) and $\nu$ are
then completely characterised by their Markov matrices.

\subsection{Stochastic transformations}

\begin{definition}[Markov Kernel / Stochastic Transformation]
A \textbf{stochastic transformation} is a function $Q: X \times \mathscr{F}
\to [0,1]$ such that, for fixed $E \in \mathscr{F}$, $Q(\cdot, E)$ is
measurable, and for fixed $x \in X$, $Q(x, \cdot)$ is a probability measure.
\end{definition}

In the discrete case $Q$ is an $M \times K$ Markov matrix with
$q_{lk} \geq 0$ and $\sum_{k} q_{lk} = 1$ for every row $l$.

\begin{definition}[Sufficiency]
$\mu$ is \textbf{sufficient for} $\nu$ if there exists a stochastic
transformation $Q$ such that
\[
\nu(\mu) = Q\,\mu,
\]
meaning the experiment $\nu$ can be replicated from $\mu$ by passing the
signal $\tilde{x}_\mu$ through the randomisation $Q$.
\end{definition}

Sufficiency captures the idea that an observer of $\tilde{x}_\mu$ can
always \emph{generate} the distribution of $\tilde{x}_\nu$ by a further
randomisation, but not vice versa.

%----------------------------------------------------------------------
\section{Three Equivalent Criteria for ``More Informative''}
%----------------------------------------------------------------------

\subsection{The economic criterion (Bonnenblust--Shapley--Sherman)}

Let $A$ be a compact convex set of actions and $u: A \times S \to \mathbb{R}$
a utility function.  An information user observes a signal $x \in X$, updates
beliefs about $\tilde{s}$, and chooses an action $d(x) \in A$ to maximise
expected utility.

\medskip
\noindent\textbf{Posterior vectors.}  Let $p = (p_1, \ldots, p_N)$ be the
\emph{a priori} distribution over $S$ with $\sum_i p_i = 1$ and $p_i \geq 0$.
Define the simplex
\[
P = \bigl\{(p_1, \ldots, p_N) : p_i \geq 0,\; \textstyle\sum_i p_i = 1\bigr\}.
\]
When experiment $\mu$ is conducted and $\tilde{x}_\mu = x$ is observed, Bayes'
rule gives the posterior
\[
\Pr(\tilde{s} = s_i \mid \tilde{x}_\mu = x)
= \frac{\Pr(\tilde{x}_\mu = x \mid \tilde{s} = s_i)\, p_i}
       {\sum_j \Pr(\tilde{x}_\mu = x \mid \tilde{s} = s_j)\, p_j}.
\]
The experiment $\mu$ therefore maps $P$ into a \emph{random posterior}
$p^{\mu}(x) \in P$ whose expectation equals the prior:
$\mathbb{E}[p^{\mu}] = p$.

\medskip
\noindent\textbf{The value of information.}  The set of possible expected
utility vectors achievable under experiment $\mu$ given prior $p$ is
\[
B(\mu, A) = \bigl\{v \in \mathbb{R}^N :
  v = \bigl(v_1(f), \ldots, v_N(f)\bigr) \text{ for some }
  \mathscr{F}\text{-measurable } f: X \to A \bigr\},
\]
where $v_i(f) = \int_X u(f(x), s_i)\,\mu_i(dx)$.
The highest expected utility available to any user who observes $\tilde{x}_\mu$
is $\max_{v \in B(\mu,A)} c \cdot v$ for any direction $c \in \mathbb{R}^N$.

\begin{definition}[Economic Criterion]
$\mu$ is at least as informative as $\nu$ (written $\mu \geq \nu$) in the sense
of Bonnenblust--Shapley--Sherman if, for every compact convex $A$ and every
$p \in P$,
\[
B(\mu, A) \supseteq B(\nu, A),
\]
i.e., every expected-utility vector achievable under $\nu$ is also achievable
under $\mu$.
\end{definition}

Equivalently, every potential information user prefers to observe $\tilde{x}_\mu$
rather than $\tilde{x}_\nu$.

\subsection{The sufficiency criterion (Blackwell)}

\begin{definition}[Blackwell Sufficiency Criterion]
$\mu \geq \nu$ in Blackwell's sense if there exists a stochastic transformation
$Q: X \to Y$ such that $\nu = Q \circ \mu$; that is,
\[
\nu_i(E) = Q(\mu_i)(E)
\quad \forall\, E \in \mathscr{G},\; i = 1, \ldots, N,
\]
where $\mathscr{G}$ is the $\sigma$-algebra on $Y$.
\end{definition}

\subsection{The uncertainty-function criterion (DeGroot)}

\citet{degroot1962} defines an \textbf{uncertainty function} $U: P \to \mathbb{R}$
to be any concave function on the simplex $P$.  The prototypical example is the
Shannon entropy
\[
U(p) = -\sum_{i=1}^{N} p_i \log p_i.
\]

\begin{definition}[DeGroot Uncertainty Criterion]
$\mu$ \textbf{reduces expected uncertainty more than} $\nu$ if, for every concave
$U: P \to \mathbb{R}$,
\[
\int_P U(p)\,\hat\mu^c(dp)
\;\leq\;
\int_P U(p)\,\hat\nu^c(dp),
\]
where $\hat\mu^c$ is the distribution of the posterior $p^\mu$ induced by
experiment $\mu$ given prior $c = (1/N, \ldots, 1/N)$.
\end{definition}

Because $U$ is concave, Jensen's inequality guarantees
$\int U(p)\,\hat\mu^c(dp) \leq U(p)$ for any experiment, so conducting an
experiment weakly reduces expected uncertainty.  The criterion asks whether
$\mu$ \emph{always} reduces it at least as much as $\nu$.

%----------------------------------------------------------------------
\section{The Main Theorem}
%----------------------------------------------------------------------

\begin{theorem}[Blackwell, 1951b; Bonnenblust et al., 1949; DeGroot, 1962]
\label{thm:blackwell}
The following three conditions are equivalent:
\begin{enumerate}[label=(\roman*)]
\item $B(\mu, A) \supseteq B(\nu, A)$ for every compact convex action set $A$
      and every prior $p \in P$ \emph{(economic criterion)};
\item there exists a stochastic transformation $Q$ from $X$ to $Y$ such that
      $\nu = Q \circ \mu$ \emph{(sufficiency criterion)};
\item $\int_P U(p)\,\hat\mu^c(dp) \leq \int_P U(p)\,\hat\nu^c(dp)$ for every
      concave $U$ and the uniform prior $c$ \emph{(uncertainty criterion)}.
\end{enumerate}
\end{theorem}

\noindent Kihlstrom's paper provides a unified Bayesian proof of all three
equivalences; we sketch the key ideas below.

%----------------------------------------------------------------------
\section{Bayesian Interpretation: Posterior Distributions and Standard Experiments}
%----------------------------------------------------------------------

\subsection{The posterior as a sufficient statistic}

A central object in Kihlstrom's analysis is the \textbf{standard experiment}
${}^c\mu^*$ associated with experiment $\mu$ and prior $c$.  When signal $x$
is observed, form the posterior
\[
p^\mu(x) = \bigl(\Pr(s_1 \mid x),\, \ldots,\, \Pr(s_N \mid x)\bigr) \in P.
\]
The standard experiment simply records this posterior: it maps the prior $c$
to the random variable $p^\mu(x) \in P$.  Its distribution $\hat\mu^c$ on $P$
satisfies the \emph{mean-preservation} property
\[
\int_P p\;\hat\mu^c(dp) = c.
\]
Two experiments that produce the same distribution of posteriors are
informationally equivalent.  The standard experiment is the minimal sufficient
statistic, capturing everything the data reveal about $\tilde{s}$.

\subsection{The Blackwell--Bayesian comparison}

Using standard-experiment language, \citet{kihlstrom1984} restates the key
comparison as follows.  Let $\hat\mu$ and $\hat\nu$ denote the distributions
of $p^\mu$ and $p^\nu$ on $P$ (the distributions of posterior vectors under
experiments $\mu$ and $\nu$ respectively).

\begin{theorem}[Kihlstrom's Reformulation]
$\mu \geq \nu$ if and only if $\hat\mu$ is a \textbf{mean-preserving spread}
of $\hat\nu$; that is, $\hat\mu$ second-order stochastically dominates
$\hat\nu$ in the sense that
\[
\int_P g(p)\,\hat\mu(dp) \;\geq\; \int_P g(p)\,\hat\nu(dp)
\]
for every convex function $g: P \to \mathbb{R}$.
\end{theorem}

The intuition is direct: a more informative experiment pushes posterior
beliefs further from the prior on average.  Posteriors are more dispersed,
reflecting a greater ``resolution'' of uncertainty.  Because every concave
$U$ is the negative of a convex function, the uncertainty criterion (iii) of
Theorem~\ref{thm:blackwell} is immediately implied.

\subsection{The stochastic transformation as a mean-preserving randomisation}

Kihlstrom proves the equivalence (i)$\Leftrightarrow$(ii) by explicit
construction.  Given $\mu \geq \nu$ in the economic sense, define $p^0 = p$
(the prior).  The key observation is that
\[
p_i^0 = \Pr(\tilde{s} = s_i) = \int_Y \Pr(\tilde{s} = s_i \mid \tilde{x}_\nu = y)\,
\Pr(\tilde{x}_\nu \in dy)
\]
and likewise for $\mu$.  Because $\mu$ achieves at least the value of $\nu$
for every user, there must exist a stochastic transformation $D(p^0, \cdot)$
on $P$ such that observing $\tilde{x}_\mu$ has the same distribution as
observing $p^\nu$ and then applying $D$.  Formally, $D$ is mean-preserving:
\[
\int_P p\,D(p^0, dp) = p^0.
\]
Setting $Q = D$ explicitly constructs the required Markov kernel, establishing
(ii).

%----------------------------------------------------------------------
\section{DeGroot's Uncertainty Function: Extensions}
%----------------------------------------------------------------------

\subsection{Concavity and the value of information}

\citet{degroot1962} formalises the notion that information has value by
introducing a general \textbf{uncertainty function} $U: P \to \mathbb{R}$.
He requires $U$ to be:
\begin{itemize}
\item \emph{Concave}: by Jensen's inequality, observing any signal weakly
      reduces expected uncertainty, $\mathbb{E}[U(p^\mu)] \leq U(p)$.
\item \emph{Symmetric}: $U$ depends on the components of $p$ but not their
      labelling.
\item \emph{Normalised}: $U$ attains its maximum at $p = (1/N, \ldots, 1/N)$
      (maximum uncertainty) and its minimum at the vertices of $P$ (certainty).
\end{itemize}
The \textbf{value of experiment $\mu$ given prior $p$} is then
\[
I(\tilde{x}_\mu;\, \tilde{s};\, U)
= U(p) - \mathbb{E}[U(p^\mu)],
\]
the expected reduction in $U$.  A key property proved by DeGroot and invoked
by Kihlstrom is that $\mu \geq \nu$ \emph{iff} $I(\tilde{x}_\mu; \tilde{s}; U)
\geq I(\tilde{x}_\nu; \tilde{s}; U)$ for every concave $U$ (Theorem 3 in
Kihlstrom).

\subsection{Entropy as a special case}

\citet{kihlstrom1984} notes that the canonical example of an uncertainty function
is Shannon entropy,
\[
U(p) = -\sum_{i=1}^{N} p_i \log p_i.
\]
Kihlstrom does not expand this special case further in his exposition.  One can
verify, however, that when $U$ is entropy and the prior is the uniform
$c = (1/N, \ldots, 1/N)$, DeGroot's value-of-information formula becomes
\[
I(\tilde{x}_\mu, c;\, U)
= U(c) - \int_P U(p)\,\hat\mu^c(dp)
= \log N - \int_P \Bigl(-{\textstyle\sum_i} p_i \log p_i\Bigr)\hat\mu^c(dp),
\]
since $U(c) = -\sum_i \tfrac{1}{N}\log\tfrac{1}{N} = \log N$.  The integral
$\int_P \bigl(-\sum_i p_i \log p_i\bigr)\hat\mu^c(dp)$ equals the expected
posterior entropy $H(\tilde{s} \mid \tilde{x}_\mu)$, so the expression reduces
to $\log N - H(\tilde{s} \mid \tilde{x}_\mu)$, the information gain
(mutual information between $\tilde{x}_\mu$ and $\tilde{s}$ under the uniform
prior).  This identification is a consequence of the framework, not something
stated explicitly by Kihlstrom.

The Blackwell ordering implies $I(\tilde{x}_\mu, c; U) \geq I(\tilde{x}_\nu, c; U)$
when $U$ is entropy.  The converse fails: entropy-valued information gain alone
does not determine the full Blackwell ordering, since the ordering requires the
inequality for \emph{every} concave $U$.

\subsection{Connection to second-order stochastic dominance}

The uncertainty-function representation makes explicit the connection between
Blackwell's theorem and second-order stochastic dominance (SOSD) of posterior
distributions.  Because $U$ is concave, $-U$ is convex, and the condition
\[
\mathbb{E}[U(p^\mu)] \;\leq\; \mathbb{E}[U(p^\nu)]
\quad \text{for all concave } U
\]
is precisely the statement that $\hat\mu$ dominates $\hat\nu$ in the SOSD
sense on $P$.  The Blackwell ordering on experiments is therefore isomorphic
to the SOSD ordering on distributions of posteriors.

%----------------------------------------------------------------------
\section{The Proof Architecture}
%----------------------------------------------------------------------

Kihlstrom's proof establishes the chain
\[
\text{(i) Economic} \;\Longleftrightarrow\; \text{(ii) Sufficiency}
\;\Longleftrightarrow\; \text{(iii) Uncertainty}.
\]

\paragraph{(ii) $\Rightarrow$ (i).}
If $Q \circ \mu = \nu$, any decision rule $f$ based on $\tilde{x}_\nu$ can
be replicated by randomising: first observe $\tilde{x}_\mu$, then draw
$\tilde{x}_\nu \sim Q(\tilde{x}_\mu, \cdot)$, and apply $f$.  Hence
$B(\nu, A) \subseteq B(\mu, A)$.

\paragraph{(i) $\Rightarrow$ (ii).}
This is the difficult direction.  Kihlstrom uses a separating-hyperplane
argument on the set of achievable utility vectors.  Since
$B(\mu, A) \supseteq B(\nu, A)$ for all compact convex $A$, standard duality
(``for every $c \in \mathbb{R}^N$, $\max_{v \in B(\mu,A)} c \cdot v \geq
\max_{v \in B(\nu,A)} c \cdot v$'') implies the existence of a stochastic
transformation $D$ mapping posteriors of $\nu$ to posteriors of $\mu$ in a
mean-preserving way.  Constructing $Q$ from $D$ completes the argument.

\paragraph{(ii) $\Leftrightarrow$ (iii).}
Given $Q$, Jensen's inequality applied to any concave $U$ directly gives
$\mathbb{E}[U(p^\mu)] \leq \mathbb{E}[U(p^\nu)]$.  Conversely, \citet{blackwell1953}
proves that the condition on all concave $U$ forces the existence of $Q$.

\paragraph{Standard experiments and the proof.}
A key simplification in Kihlstrom's exposition is the use of the standard
experiment ${}^c\mu^*$: since ${}^c\mu^*$ is a sufficient statistic for $\mu$,
it suffices to verify the conditions for standard experiments.  Theorem 5 in
\citet{kihlstrom1984} asserts that $\mu^* \geq \nu^*$ if and only if there
exists a mean-preserving stochastic transformation $D$ from $P \times \mathscr{P}$
into $P$ such that, for almost all $(p^0, R)$, the random variable
$\tilde{p}_\mu$ (posterior under $\mu$) and $\tilde{p}_\nu$ (posterior under $\nu$)
have the same joint distribution as $(\tilde{p}_\mu, D(\tilde{p}_\nu, \cdot))$.
The fact that $D$ is mean-preserving—i.e., $D(p^0, \cdot)$ has mean $p^0$
for all $p^0$—guarantees no information is added by the randomisation, only destroyed.

%----------------------------------------------------------------------
\section{Applications}
%----------------------------------------------------------------------

\subsection{Product quality information (Kihlstrom, 1974)}

\citet{kihlstrom1974a} applies Blackwell's theorem to consumer demand for
information about product quality.  The unknown state $\tilde{s}$ is a
product parameter $\theta$.  A consumer can observe $\theta$ units of
information at cost proportional to $\theta$.  As $\theta$ rises, the
experiment becomes more informative in the Blackwell sense.  The consumer's
demand for information is determined by weighing the marginal utility of the
standard experiment against its cost.  The Blackwell ordering certifies that
``more information is always better off'' for every expected-utility maximiser,
provided information is free.

\subsection{Sequential experimental design (DeGroot, 1962)}

\citet{degroot1962} applies the uncertainty function framework to sequential
design.  In this context $\tilde{x}_\theta$ is a normal random variable with
mean $\theta$ and variance $1/\theta$.  The standard experiment reduces to
observing sample averages.  The choice criterion is the minimisation of
expected uncertainty across all $n$ experiments.  \citeauthor{degroot1962}
proves, using Blackwell's Theorem, that the optimal sequential strategy is to
use the experiment $E$ that minimises expected uncertainty regardless of the
outcomes of earlier experiments.

%----------------------------------------------------------------------
\section{Summary}
%----------------------------------------------------------------------

Blackwell's theorem identifies a \emph{partial order} on statistical
experiments that has three equivalent characterisations: economic dominance
(every decision maker is at least as well off), sufficiency (the more
informative signal can replicate the less informative one by randomisation),
and uncertainty reduction (every concave measure of uncertainty is reduced
more in expectation).

Kihlstrom's Bayesian exposition clarifies the structure of the theorem
by placing the posterior distribution at the centre.  The more informative
experiment generates a distribution of posterior beliefs that is a
mean-preserving spread of the less informative one.  This geometric insight
connects Blackwell's ordering to second-order stochastic dominance and gives
the uncertainty-function criterion a transparent interpretation: because $U$
is concave, mean-preserving spreads of posteriors correspond to higher expected
$U$, i.e., lower expected uncertainty.

DeGroot's extension broadens the criterion from specific utility functions to
the entire class of concave uncertainty functions, reinforcing the generality
of the result.  Together, these formulations provide an analytically powerful
toolkit for ranking information structures in macroeconomics, mechanism design,
and statistical decision theory.

\bigskip
\bibliographystyle{plainnat}
\bibliography{blackwell_kihlstrom}

\end{document}
